\documentclass[11pt]{article}
\usepackage[margin=1in]{geometry}
\usepackage{amsmath,amssymb,amsthm}
\usepackage{graphicx}
\usepackage[hidelinks]{hyperref}

\newtheorem{theorem}{Theorem}
\newtheorem{lemma}{Lemma}
\newtheorem{remark}{Remark}

\newcommand{\U}{\mathcal U}
\newcommand{\dens}{\operatorname{dens}}
\newcommand{\fin}{\operatorname{fin}}
\newcommand{\norm}[1]{\left\lVert #1 \right\rVert}
\newcommand{\Linfty}{\mathcal L_{\infty}}

\graphicspath{{runs/fa_banach_001/solutions/full/1406.6733_l_infty_ultraproducts_aleph_injective/}}

\title{\(\mathcal L_\infty\) Ultraproducts Over Good Ultrafilters Are \(\aleph\)-Injective}
\author{FA Banach run \texttt{fa\_banach\_001}, agent \texttt{agent\_lane\_06}}
\date{2026-06-19}

\begin{document}
\maketitle

\section*{Status}

This packet records a candidate full solution, likely valid, to Open Problem
(6) of Aviles, Cabello Sanchez, Castillo, Gonzalez, and Moreno \cite{ACCGM}:
whether every ultraproduct over a countably incomplete \(\aleph\)-good
ultrafilter is \(\aleph\)-injective as long as it is an
\(\mathcal L_\infty\)-space.

\section*{Source Passage}

\begin{center}
\includegraphics[width=0.98\linewidth]{figures/open_problem_crop.png}
\end{center}

\section*{Transcription of the Open Problem}

Open Problem (6) asks:
\begin{quote}
``Prove or disprove that every ultraproduct built over a countably incomplete,
\(\aleph\)-good ultrafilter is \(\aleph\)-injective as long as it is a
\(\mathcal L_\infty\)-space.''
\end{quote}

\section*{Definitions Used}

Let \(\aleph\) be an uncountable cardinal.  A Banach space \(E\) is
\(\aleph\)-injective if, whenever \(Y\subset X\), \(\dens X<\aleph\), and
\(t:Y\to E\) is an operator, there is an extension \(T:X\to E\).  It is
\((\lambda,\aleph)\)-injective if \(T\) can always be chosen with
\(\norm T\le \lambda\norm t\).

Following the source paper, an ultrafilter \(\U\) is countably incomplete if
there are \(I_n\in\U\) with \(I_{n+1}\subset I_n\) and
\(\bigcap_n I_n=\varnothing\).  It is \(\aleph\)-good if every monotone map
from the finite subsets of any set of cardinality \(<\aleph\) into \(\U\) has
a multiplicative refinement.

A Banach space \(E\) is an \(\mathcal L_{\infty,\lambda}\)-space if every
finite-dimensional subspace of \(E\) is contained in a finite-dimensional
subspace whose Banach--Mazur distance to some \(\ell_\infty^n\) is at most
\(\lambda\).  An \(\mathcal L_\infty\)-space means
\(\mathcal L_{\infty,\lambda}\) for some finite \(\lambda\).

\section*{Proof Intuition}

The source theorem proves the isometric case by turning a family of fewer than
\(\aleph\) balls into fewer than \(\aleph\) internal sets and applying
good-ultrafilter saturation.  For a general \(\mathcal L_\infty\) ultraproduct,
one should not ask for exact ball intersections.  The right finite-level
input is instead the standard local extension property of
\(\mathcal L_{\infty,\lambda}\)-spaces: finite-dimensional extension problems
can be solved with constant \(\lambda\).

The good ultrafilter then upgrades finite solvability to solvability for all
constraints indexed by a dense set of size \(<\aleph\).  We encode a candidate
extension as a function from a fixed rational dense subspace into the
ultraproduct.  Linearity, agreement with the original operator, and the
\(\lambda\)-norm bound are all closed finite constraints.  Every finite
subcollection is solved by the finite-dimensional extension lemma, so
saturation realizes the whole type.

\section*{Theorem}

\begin{theorem}
Let \(\U\) be a countably incomplete \(\aleph\)-good ultrafilter, and let
\[
  E=[E_i]_{\U}
\]
be an ultraproduct of Banach spaces.  If \(E\) is an
\(\mathcal L_{\infty,\lambda}\)-space, then \(E\) is
\((\lambda,\aleph)\)-injective.  Consequently, if \(E\) is an
\(\mathcal L_\infty\)-space, then \(E\) is \(\aleph\)-injective.
\end{theorem}

\begin{lemma}[Finite-dimensional extension]\label{lem:finite}
Let \(E\) be an \(\mathcal L_{\infty,\lambda}\)-space.  If
\(A\subset B\) are finite-dimensional Banach spaces and \(u:A\to E\) is an
operator, then \(u\) has an extension \(\widetilde u:B\to E\) with
\(\norm{\widetilde u}\le \lambda\norm u\).
\end{lemma}

\begin{proof}
Let \(F=u(A)\).  By the \(\mathcal L_{\infty,\lambda}\) property, \(F\) is
contained in a finite-dimensional subspace \(G\subset E\) and there is an
isomorphism \(S:G\to \ell_\infty^n\) such that
\(\norm S\norm{S^{-1}}\le\lambda\).  Since \(\ell_\infty^n\) is 1-injective,
\(Su:A\to\ell_\infty^n\) extends to an operator \(v:B\to\ell_\infty^n\) with
\(\norm v=\norm{Su}\).  Then \(S^{-1}v:B\to G\subset E\) extends \(u\), and
\[
  \norm{S^{-1}v}
  \le \norm{S^{-1}}\norm S\norm u
  \le \lambda\norm u.
\]
\end{proof}

\section*{Proof of the Theorem}

We use the saturation fact stated in the proof of Theorem 5.1 of
\cite{ACCGM}: if \(\U\) is countably incomplete and \(\aleph\)-good, then any
family of fewer than \(\aleph\) internal subsets of a set-theoretic
ultraproduct with the finite intersection property has nonempty intersection.

Let \(Y\subset X\) be Banach spaces with \(\dens X<\aleph\), and let
\(t:Y\to E\) be an operator.  The case \(t=0\) is trivial, so put
\(M=\norm t>0\).  Choose a rational vector subspace \(D\subset X\) which is
dense in \(X\), has cardinality \(<\aleph\), and such that
\[
  D_Y:=D\cap Y
\]
is dense in \(Y\).  For each \(y\in D_Y\), fix a representative
\((t_i(y))_{i\in I}\) of \(t(y)\in [E_i]_\U\).

Let \(S_i\) be the set of all functions \(\varphi_i:D\to E_i\).  An element
\(\langle \varphi_i\rangle_\U\) of the set-theoretic ultraproduct
\(\langle S_i\rangle_\U\) determines a candidate map
\[
  T_0:D\to E,\qquad T_0(x)=[(\varphi_i(x))]_\U .
\]
We now impose fewer than \(\aleph\) internal constraints on
\(\langle \varphi_i\rangle_\U\).

First, for \(x,x'\in D\), \(q\in\mathbb Q\), and \(m\in\mathbb N\), impose the
approximate linearity constraints
\[
  \norm{\varphi_i(x+x')-\varphi_i(x)-\varphi_i(x')}\le 1/m,
  \qquad
  \norm{\varphi_i(qx)-q\varphi_i(x)}\le 1/m.
\]
Second, for \(y\in D_Y\) and \(m\in\mathbb N\), impose the agreement constraint
\[
  \norm{\varphi_i(y)-t_i(y)}\le 1/m.
\]
Finally, for every finite tuple \(x_1,\ldots,x_n\in D\), rational scalars
\(q_1,\ldots,q_n\), and \(m\in\mathbb N\), impose the norm constraint
\[
  \norm{\sum_{j=1}^n q_j\varphi_i(x_j)}
  \le
  \lambda M \norm{\sum_{j=1}^n q_j x_j}+1/m.
\]
Each displayed condition defines an internal subset of
\(\langle S_i\rangle_\U\), and the total number of such constraints is
\(<\aleph\).

It remains to check finite satisfiability.  Take finitely many of the above
constraints.  Only finitely many vectors of \(D\) and \(D_Y\) occur.  Let
\[
  B=\operatorname{span}\{ \text{all occurring vectors of }D\},\qquad
  A=B\cap Y .
\]
Since \(B\) is finite-dimensional, Lemma \ref{lem:finite} extends
\(t|_A:A\to E\) to an operator \(u:B\to E\) with
\(\norm u\le \lambda M\).  Choose representatives \((u_i(x))_{i\in I}\) for
the finitely many values \(u(x)\) attached to vectors that occur in the chosen
constraints, and define \(\varphi_i\) on those relevant vectors by these
representatives, extending it arbitrarily to all of \(D\).  Since the required
equalities and
norm bounds hold in the ultraproduct, the set of \(i\)'s on which the finite
list of approximate coordinate inequalities holds belongs to \(\U\).  Hence
the corresponding finite intersection of internal sets is nonempty.

By the good-ultrafilter saturation theorem, there is
\(\langle \varphi_i\rangle_\U\in \langle S_i\rangle_\U\) satisfying all
constraints.  The associated map \(T_0:D\to E\) is exactly
\(\mathbb Q\)-linear, agrees with \(t\) on \(D_Y\), and satisfies
\[
  \norm{T_0(x)}\le \lambda M\norm x\qquad(x\in D).
\]
Thus \(T_0\) extends uniquely by continuity to a bounded linear operator
\[
  T:X\to E
\]
with \(\norm T\le\lambda\norm t\).  Since \(D_Y\) is dense in \(Y\) and
\(T_0=t\) on \(D_Y\), the extension satisfies \(T|_Y=t\).

Therefore \(E\) is \((\lambda,\aleph)\)-injective.  If \(E\) is merely an
\(\mathcal L_\infty\)-space, then \(E\) is
\(\mathcal L_{\infty,\lambda}\) for some finite \(\lambda\), so the preceding
paragraph gives \(\aleph\)-injectivity.  This answers Open Problem (6)
affirmatively.
\(\square\)

\section*{Verification Notes}

No numerical verification is involved.  The finite-dimensional extension lemma
is the standard \(\mathcal L_{\infty,\lambda}\) argument through
\(\ell_\infty^n\).  The only set-theoretic input is precisely the internal-set
saturation property quoted and used in the source paper's proof of Theorem
5.1.

The result does not address Open Problem (5), which asks for universal
\(\aleph\)-injectivity of Lindenstrauss ultraproducts in ZFC.  The proof here
requires \(\dens X<\aleph\), exactly as in the definition of ordinary
\(\aleph\)-injectivity.

\section*{Novelty and Literature Check}

The local run indexes were searched for arXiv:1406.6733, ``aleph-injective'',
``aleph-good'', ``good ultrafilter'', ``countably incomplete'',
``\(\mathcal L_\infty\) ultraproduct'', and related phrases.  No existing
packet or attempt covers Open Problem (6).

External web searches on 2026-06-19 for exact phrases around
\(\aleph\)-good ultrafilters, \(\mathcal L_\infty\) ultraproducts, and
\(\aleph\)-injectivity found the source paper \cite{ACCGM} and the earlier
separably-injective paper \cite{ACCGMsep}, but no later paper explicitly
resolving this higher-cardinal problem.

Human review should focus on the saturation coding in the proof: the ambient
set-theoretic ultraproduct is built from the sets \(S_i\) of all functions
\(D\to E_i\), and each linearity, agreement, or norm bound is imposed as an
internal coordinatewise condition with \(1/m\) tolerance.

\begin{thebibliography}{9}

\bibitem{ACCGM}
A. Aviles, F. Cabello Sanchez, J. M. F. Castillo, M. Gonzalez, and Y. Moreno,
\emph{\(\aleph\)-injective Banach spaces and \(\aleph\)-projective compacta},
arXiv:1406.6733.

\bibitem{ACCGMsep}
A. Aviles, F. Cabello Sanchez, J. M. F. Castillo, M. Gonzalez, and Y. Moreno,
\emph{On separably injective Banach spaces and Corrigendum to ``On separably
injective Banach spaces''},
arXiv:1103.6064.

\bibitem{LT}
J. Lindenstrauss and L. Tzafriri,
\emph{Classical Banach Spaces I and II},
Classics in Mathematics, Springer, 1996.

\end{thebibliography}

\end{document}
